\chapter{Conclusions}
\label{chap:conclusions}

This thesis presented the design, implementation, and experimental use of a
lightweight framework for the generation and management of Kathar\'a network
laboratories. The work originated from a practical problem: even relatively
small emulated networks require several artefacts to remain consistent with one
another. A topology must define the correct collision domains, every node must
receive the intended interfaces and IP addresses, gateways and static routes
must match the addressing plan, routing daemons must be enabled and configured,
and application-specific files must be placed in the correct node directories.
When these tasks are performed manually, the resulting laboratory is often
correct only after several iterations of editing and troubleshooting. The
framework developed in this thesis addresses this problem by moving the primary
description of a laboratory to a declarative YAML configuration and deriving the
Kathar\'a artefacts from that source.

The contribution is deliberately modest in scope. The project does not attempt
to replace Kathar\'a, Docker, or FRRouting, and it does not introduce a new
network-emulation runtime; instead, it preserves the standard Kathar\'a
workflow described by the project documentation
\cite{kathara_website,kathara_manpage}. It adds a small abstraction layer that
makes a selected class of laboratories easier to describe, validate, regenerate,
and reuse while preserving compatibility with the normal Kathar\'a workflow.
This design choice shaped the complete implementation and is also central to the
assessment presented in this concluding chapter.

\section{Summary of Contributions}
\label{sec:conclusion-contributions}

The first contribution of the work is the definition of a declarative
configuration model for Kathar\'a laboratories. A single YAML document can
represent the laboratory name and metadata, node roles, interfaces, collision
domains, IP addressing, default gateways, static routes, startup commands, and
routing-related information. The same model supports several node categories
and can express both simple host-router-server laboratories and more structured
routed environments. The use of YAML keeps the source configuration readable
and versionable, while the generated Kathar\'a directory remains a derived
artefact rather than the authoritative description of the topology.

A second contribution is the validation stage introduced before generation.
The framework first checks the minimum structural requirements of the document
and then performs semantic checks on the network description. These checks
include node and collision-domain naming, node types, IP syntax, duplicate
addresses, use of network or broadcast addresses, inconsistent or overlapping
subnets, default gateways that are not locally reachable, routing-device
consistency, router identifiers, OSPF network declarations, and BGP parameters
and neighbour relationships. The distinction between blocking errors and
warnings allows the framework to reject configurations that cannot be generated
safely while still reporting suspicious situations that may be intentional.
The \cliopt{validate-only} mode also makes validation useful independently of
generation.

The third contribution is the automatic production of standard Kathar\'a
artefacts. Starting from a validated configuration, the generator creates the
laboratory directory, writes \texttt{lab.conf}, produces the startup files for
individual nodes, selects suitable container images, and creates the FRRouting
files required by routing-capable devices. OSPF, RIP, and BGP are supported by
the current generation logic, while the configuration model also provides an
escape mechanism for cases in which a complete custom FRRouting configuration
is more appropriate than the declarative subset implemented by the framework.
Node-specific directories can additionally be imported into the generated
laboratory, allowing application scripts and configuration files to coexist
with the automatically generated network artefacts.

The generator also provides optional Wireshark integration. The framework can
identify the collision domains defined by the topology and attach a dedicated
capture container to selected networks. This avoids requiring separate Docker
or packet-capture setup after generation and keeps the observation point part
of the laboratory itself. Together with the lifecycle wrappers for starting,
stopping, and removing laboratories, these features create a complete workflow
from configuration to execution while still relying on standard Kathar\'a
commands underneath.

A further contribution is the explicit separation between configuration,
generation, execution, and experimentation. The architecture described in
\Cref{chap:requirements-architecture} is reflected by the implementation: YAML
files describe the desired laboratory, the loader and validator analyse that
description, the generator produces the runtime artefacts, and separate tools
operate on the generated environment. This separation prevents the topology
description from becoming entangled with one specific experiment. It also makes
it possible to regenerate a laboratory after changing the configuration instead
of applying persistent modifications directly to generated files.

Finally, the framework was exercised in a non-trivial experimental setting
rather than being evaluated only through minimal examples. The case study in
\Cref{chap:experimental-evaluation} combined eight legitimate clients, multiple
routers, OSPF-based connectivity, an application service, concurrent
measurements, server-side metrics, and a bounded adverse workload. The purpose
of that scenario was not to make attack generation the central contribution of
the thesis. It was used to verify that laboratories produced by the framework
could support coordinated, repeatable measurements in substantially different
operating conditions. A second topology placed a rate-limiting reverse proxy in
front of the service, demonstrating that the same measurement schema can also
support an explicit mitigation comparison.

From a software-engineering perspective, these contributions are also the
result of successive refinement rather than independent features selected in
advance. The project moved from basic topology generation to pre-generation
validation, richer routing support, diagnostic and lifecycle facilities,
application-level testing, concurrent measurements, controlled experimentation,
and finally robustness and scalability analysis. Several final behaviours are
direct consequences of problems observed during earlier iterations: service
checks were separated from reachability after runtime testing showed that a
reachable node could still provide no usable application response; concurrent
rounds replaced sequential client execution when temporal comparability became
necessary; readiness checks were strengthened after an early baseline was
contaminated by unavailable routes; and the final stress tests exposed the
performance cost of defensive duplicated validation. The implemented
architecture therefore records both successful design decisions and trade-offs
identified through verification.

These contributions can be summarised as a progression from a declarative
network description to an executable and observable environment:

\begin{center}
\begin{tikzpicture}[
    node distance=1.25cm,
    every node/.style={font=\small},
    box/.style={draw, rounded corners, minimum width=3.25cm, minimum height=0.72cm, align=center},
    arrow/.style={-{Latex[length=2mm]}, thick}
]
\node[box] (yaml) {Declarative YAML\\configuration};
\node[box, below=of yaml] (validate) {Structural and\\semantic validation};
\node[box, below=of validate] (generate) {Kathar\'a artefact\\generation};
\node[box, below=of generate] (runtime) {Laboratory execution\\and inspection};
\node[box, below=of runtime] (experiment) {Repeatable measurement\\and experimentation};
\draw[arrow] (yaml) -- (validate);
\draw[arrow] (validate) -- (generate);
\draw[arrow] (generate) -- (runtime);
\draw[arrow] (runtime) -- (experiment);
\end{tikzpicture}
\end{center}

The resulting framework remains intentionally small, but the complete chain is
important. Automation is useful only if the input can be checked, the output is
understandable, and the generated environment remains compatible with the tools
that users already expect to use. For this reason, the project favours explicit
files and standard Kathar\'a conventions over a more opaque orchestration layer.

\section{Main Findings and Overall Assessment}
\label{sec:conclusion-findings}

The implementation confirms that a declarative front end can reduce the amount
of repetitive work required to construct Kathar\'a laboratories without hiding
the resulting network configuration. The generated \texttt{lab.conf}, startup
files, and FRRouting configuration remain directly inspectable. This is
particularly useful in an educational or experimental context, where the
configuration itself is often as important as the running network. The
framework therefore automates file construction while preserving the ability to
connect to a node, inspect interfaces and routes, use \texttt{vtysh}, capture
packets, and modify or troubleshoot the environment with ordinary Kathar\'a and
Linux tools.

The validation layer proved equally important. Several classes of mistakes can
be identified before any container is started, including addressing conflicts,
gateways outside directly connected networks, and inconsistent routing
parameters. Detecting these conditions at configuration time does not guarantee
that the running network will behave exactly as intended, but it reduces a set
of errors that would otherwise appear later as connectivity or routing
failures. This distinction between static correctness and runtime correctness is
also one of the reasons why runtime verification was identified in
\Cref{chap:future-developments} as a natural direction for future work.

The framework also demonstrated that the same generation process can be used
for topologies of different complexity. Beyond the example laboratories, the
translator was stress-tested in Chapter~\ref{chap:robustness-security-scalability}
with synthetic configurations containing up to 10,000 nodes and with inputs
dominated by large route or command sets. These results characterise loading,
validation, and file-generation behaviour, but they do not establish that
laboratories of comparable size can be executed efficiently in Kathar\'a. The
runtime scalability of large generated laboratories therefore remains outside
the evaluated scope.

The experimental evaluation provides a more concrete assessment of the running
environment. Five complete executions were collected for each topology. Every
direct-server baseline completed all 80 legitimate application requests. In
every subsequent controlled run, 16 of 80 requests completed, corresponding to
20\% over the complete transition; all 16 occurred in the first two rounds
before workload activation. During rounds 3--10, every direct execution
therefore completed 0 of 64 legitimate requests, corresponding to 0\%
attack-active availability. At the same time, all ICMP probes and TCP
connection attempts succeeded. In each run the instrumented server reached its
configured limit of 20 active application connections and recorded 264 rejected
or \texttt{BUSY} events during the load phase.

The five reverse-proxy executions produced a different and equally stable
outcome under the same round structure. Every baseline and controlled-load
phase completed all 80 legitimate requests, so attack-active availability was
100\% in every run, compared with 0\% without the proxy. The protected backend
reached at most eight simultaneous connections and recorded no rejected or
\texttt{BUSY} event in any repetition. Across paired runs, the increase in the
per-run baseline median application response had a median of 2.167 ms and an
interquartile interval of [2.147, 2.472] ms, illustrating the
availability--overhead trade-off in the tested configuration.

The direct-server result is significant mainly because of the distinction it exposes.
The emulated routing infrastructure remained reachable and the TCP service port
continued to accept connections, yet the application was unable to complete
most legitimate requests. The experiment therefore demonstrates that ICMP
reachability, TCP connection establishment, and application availability are
separate observations. Measuring only ping latency or whether a port can be
opened would have missed the principal effect observed by the legitimate
clients. For this reason, the framework's measurement tools record several
timing and success dimensions instead of collapsing the experiment into a
single latency value.

The experimental evidence must nevertheless be interpreted within its scope.
The five paired acquisition blocks provide evidence of consistent run-to-run
behaviour for the implemented testbed under the stated parameters and clean
laboratory restarts, as detailed in \Cref{sec:repeated-run-validation}. They
remain a small sample, were collected on one host and day, and used a fixed
direct-then-mitigated order. Moreover, the mitigation result applies to the
bounded connection-oriented workload used in the laboratory and does not
establish protection against volumetric or more diverse application-layer
attacks.

This limitation is useful when assessing the framework as a research prototype.
The project succeeds in automating configuration, validation, generation, and a
repeatable experimental workflow, but it does not yet provide all of the
controls expected from a mature experimentation platform. There is no formal
schema independent of the Python implementation, no complete typed internal
network model, no automatic verification that the runtime state matches the
expected state, and no systematic benchmark of runtime execution at
large scale. These points do not invalidate the implemented workflow; they
define the boundary of the contribution and motivate the evolution path
presented in the previous chapter.

From an architectural perspective, the most important outcome is therefore not
a single generated topology or a single experimental result. It is the
separation of concerns established by the project. The source configuration
expresses intent, validation checks that intent against the supported model,
generation converts it into standard runtime artefacts, and measurement tools
observe the behaviour of the resulting network. Because each stage has a clear
responsibility, the framework can be extended without requiring the current
workflow to be discarded. Declarative experiment definitions, runtime
verification, richer reporting, topology visualisation, or additional backend
abstractions can all build on the same separation.


\subsection{Answers to the Research Questions}
\label{sec:conclusion-rq-answers}

\noindent\textbf{RQ1 --- Declarative translation.}
The framework translates the YAML description through a staged pipeline:
loading, structural and semantic validation, and generation. The output is an
ordinary Kathar\'a directory containing \texttt{lab.conf}, per-node startup
files, imported resources, and FRRouting files where required. The generated
artefacts remain inspectable and usable with standard Kathar\'a commands.

\noindent\textbf{RQ2 --- Generated configuration and deployed resources.}
Address assignments, gateways, static routes, and the supported OSPF, RIP, and
BGP forms are translated into startup or FRRouting configuration. Application
resources and startup commands can be deployed from node-specific directories,
and a Wireshark node can be attached to selected collision domains. The current
model does not, however, turn \texttt{forwarding: true} into a kernel command or
turn a \texttt{switch} role into an automatic Layer~2 bridge; such behaviour
requires explicit commands or resources.

\noindent\textbf{RQ3 --- Repeated service measurement.}
The client workflow executes concurrent rounds from several legitimate nodes
and records ICMP reachability, TCP establishment, application completion, and
timing fields in CSV form. Success requires the complete application exchange
of selection \texttt{1} followed by \texttt{pong}, so a reachable host or a
successful handshake is not misreported as an available service.

\noindent\textbf{RQ4 --- Controlled comparison and mitigation.}
The same client procedure and result schema were used for the baseline,
controlled-load, and reverse-proxy experiments. Under the tested active-load
condition, availability was 0\% in all five direct-server runs and 100\% in all
five protected runs, while backend occupancy and rejection counters supplied
independent server-side evidence. This supports the mitigation only for the
bounded workload and parameters evaluated here.

\noindent\textbf{RQ5 --- Deterministic framework and LLM assistance.}
An LLM can draft YAML, shell, or FRRouting text and help interpret results, but
plausible text is not equivalent to a checked, reproducible experiment. The
framework supplies a validated configuration interface, repeatable semantic
checks, coordinated artefact generation, and machine-readable evidence. The
two approaches are therefore complementary: flexible drafting and explanation
surround a deterministic validation--generation--measurement boundary, in line
with work combining LLM-based configuration with explicit synthesis or
verification \cite{wang2024netconfeval,lira2025llmnetcfg}.

\section{Final Remarks}
\label{sec:conclusion-final-remarks}

The objective of this thesis was to investigate whether a simple framework
could make the construction and use of Kathar\'a laboratories more systematic
without replacing the tools and file formats on which Kathar\'a is based. The
implemented prototype shows that this is feasible. A readable YAML description
can be validated and transformed into an ordinary Kathar\'a laboratory,
including node startup configuration and dynamic-routing files, while optional
resources and packet-capture support can be incorporated into the same
generation process. The resulting environment remains accessible through the
normal Kathar\'a lifecycle and can be regenerated from its declarative source.

The experimental case study further shows why this type of automation is useful.
Once the laboratory structure is reproducible, the focus of an experiment can
move from repeatedly preparing the environment to observing its behaviour. In
the evaluation performed for this thesis, concurrent measurements captured a
clear difference between normal service operation and operation under a bounded
adverse workload, while simultaneously showing that network reachability alone
was insufficient to describe application behaviour. This provides a practical
example of how the generated infrastructure can be used for networking and
security experiments rather than only for topology deployment.

The framework should therefore be considered a foundation rather than a final
platform. Its current implementation addresses a deliberately limited set of
network primitives and operational tasks, but it establishes a coherent path
from configuration to execution and measurement. The developments proposed in
\Cref{chap:future-developments}---generalised security and resilience testing,
runtime verification, improved experiment management, an optional LLM-assisted
configuration loop, and a cleaner separation between the logical model and the
execution backend---would extend this foundation while preserving its original
emphasis on explicit configuration and reproducibility.

In its present form, the project demonstrates that useful automation does not
require hiding the network from the user. By generating standard Kathar\'a
artefacts from a validated declarative source, the framework reduces repetitive
configuration while keeping the laboratory transparent, inspectable, and
compatible with existing tools. The same principle remains valid in an
LLM-assisted workflow: natural-language assistance can make configuration more
accessible, but the accepted experiment should still pass through explicit,
inspectable, and reproducible validation and generation stages. This balance
between flexible assistance, deterministic control, and visibility is the main
result of the work and provides the basis for further development in network
emulation, teaching, controlled security testing, and reproducible
experimentation.
